% =========================================================================
% Figura: Pipeline tradicional de descubrimiento y desarrollo de fármacos
% Inspirado en Chawla et al. (2024), Figura 1.1 y Hughes et al. (2011)
% Paleta de colores accesible Okabe-Ito definida en Configurations/01-Colors.sty
% =========================================================================

\begin{figure}[htbp]
\centering
\resizebox{\linewidth}{!}{%
\begin{tikzpicture}[
    font=\sffamily,
    >=Stealth,
    x=1cm, y=1cm
]

% -------------------------------------------------------------------------
% Parámetros geométricos
% -------------------------------------------------------------------------
\def\w{3.4}      % Ancho de cada bloque
\def\hArrow{1.15}% Altura de la flecha de compuestos
\def\hBadge{0.72}% Altura del distintivo de la etapa
\def\tip{0.45}   % Saliente de la punta de la flecha
\def\wing{0.16}  % Aleta de la flecha

% Coordenadas horizontales
\def\xZero{0.00}
\def\xOne{3.75}
\def\xTwo{7.50}
\def\xThree{11.25}
\def\totalW{14.65}

% =========================================================================
% FILA 1: FLECHAS CON CANTIDAD DE COMPUESTOS (EMBUDO)
% =========================================================================

% --- Etapa 1: 10.000 compuestos (Azul) ---
\filldraw[draw=figBlue, fill=figBlueTint, line width=1.1pt, rounded corners=1pt]
    (\xZero, 0) -- (\xZero + \w - \tip, 0) -- (\xZero + \w - \tip, \wing) -- (\xZero + \w, -\hArrow/2) -- (\xZero + \w - \tip, -\hArrow - \wing) -- (\xZero + \w - \tip, -\hArrow) -- (\xZero, -\hArrow) -- cycle;
\node[align=center, text=darkText] at (\xZero + \w/2 - 0.15, -\hArrow/2) {
    {\fontsize{12}{14}\selectfont \textbf{10.000}}\\[1pt]
    {\footnotesize compuestos iniciales}
};

% --- Etapa 2: 250 compuestos (Naranja) ---
\filldraw[draw=figOrange, fill=figOrangeTint, line width=1.1pt, rounded corners=1pt]
    (\xOne, 0) -- (\xOne + \w - \tip, 0) -- (\xOne + \w - \tip, \wing) -- (\xOne + \w, -\hArrow/2) -- (\xOne + \w - \tip, -\hArrow - \wing) -- (\xOne + \w - \tip, -\hArrow) -- (\xOne, -\hArrow) -- cycle;
\node[align=center, text=darkText] at (\xOne + \w/2 - 0.15, -\hArrow/2) {
    {\fontsize{12}{14}\selectfont \textbf{250}}\\[1pt]
    {\footnotesize candidatos preclínicos}
};

% --- Etapa 3: 5 compuestos (Bermellón) ---
\filldraw[draw=figVermilion, fill=figVermilionTint, line width=1.1pt, rounded corners=1pt]
    (\xTwo, 0) -- (\xTwo + \w - \tip, 0) -- (\xTwo + \w - \tip, \wing) -- (\xTwo + \w, -\hArrow/2) -- (\xTwo + \w - \tip, -\hArrow - \wing) -- (\xTwo + \w - \tip, -\hArrow) -- (\xTwo, -\hArrow) -- cycle;
\node[align=center, text=darkText] at (\xTwo + \w/2 - 0.15, -\hArrow/2) {
    {\fontsize{12}{14}\selectfont \textbf{5}}\\[1pt]
    {\footnotesize moléculas en clínica}
};

% --- Etapa 4: 1 fármaco (Verde azulado) ---
\filldraw[draw=figGreen, fill=figGreenTint, line width=1.1pt, rounded corners=1pt]
    (\xThree, 0) -- (\xThree + \w - \tip, 0) -- (\xThree + \w - \tip, \wing) -- (\xThree + \w, -\hArrow/2) -- (\xThree + \w - \tip, -\hArrow - \wing) -- (\xThree + \w - \tip, -\hArrow) -- (\xThree, -\hArrow) -- cycle;
\node[align=center, text=darkText] at (\xThree + \w/2 - 0.15, -\hArrow/2) {
    {\fontsize{12}{14}\selectfont \textbf{1}}\\[1pt]
    {\footnotesize fármaco aprobado}
};

% =========================================================================
% FILA 2: DISTINTIVOS CON EL NOMBRE DE CADA ETAPA
% =========================================================================

\node[
    draw=figBlue,
    fill=figBlueTint,
    line width=0.9pt,
    rounded corners=3pt,
    minimum width=\w cm,
    minimum height=\hBadge cm,
    text width=\w cm - 0.2cm,
    anchor=north west,
    align=center,
    inner sep=3pt
] at (\xZero, -1.35) {
    \textbf{\small\textcolor{figBlue}{Descubrimiento}}\\[1pt]
    {\footnotesize\textcolor{figBlue}{temprano}}
};

\node[
    draw=figOrange,
    fill=figOrangeTint,
    line width=0.9pt,
    rounded corners=3pt,
    minimum width=\w cm,
    minimum height=\hBadge cm,
    text width=\w cm - 0.2cm,
    anchor=north west,
    align=center,
    inner sep=3pt
] at (\xOne, -1.35) {
    \textbf{\small\textcolor{figOrange!90!black}{Fase preclínica}}\\[1pt]
    {\footnotesize\textcolor{grayText}{(\textit{in vitro} / \textit{in vivo})}}
};

\node[
    draw=figVermilion,
    fill=figVermilionTint,
    line width=0.9pt,
    rounded corners=3pt,
    minimum width=\w cm,
    minimum height=\hBadge cm,
    text width=\w cm - 0.2cm,
    anchor=north west,
    align=center,
    inner sep=3pt
] at (\xTwo, -1.35) {
    \textbf{\small\textcolor{figVermilion}{Fases clínicas}}\\[1pt]
    {\footnotesize\textcolor{figVermilion}{(Fases I--III)}}
};

\node[
    draw=figGreen,
    fill=figGreenTint,
    line width=0.9pt,
    rounded corners=3pt,
    minimum width=\w cm,
    minimum height=\hBadge cm,
    text width=\w cm - 0.2cm,
    anchor=north west,
    align=center,
    inner sep=3pt
] at (\xThree, -1.35) {
    \textbf{\small\textcolor{figGreen!80!black}{Aprobación}}\\[1pt]
    {\footnotesize\textcolor{figGreen!80!black}{regulatoria}}
};

% =========================================================================
% FILA 3: BARRAS GLOBALES DE TIEMPO Y COSTO (MÉTRICAS DEL PROCESO COMPLETO)
% Tonos neutros (slate/gris) para indicar alcance global y evitar confusión
% con los colores de las etapas específicas.
% =========================================================================
\def\hBar{0.75}  % Altura más amplia para dar respiro al texto
\def\tipBar{0.45}% Punta de la flecha
\def\wingBar{0.14}% Aleta de la flecha

% --- Barra 1: Tiempo acumulado (10 a 14 años) ---
\filldraw[draw=grayText, fill=cardBorder!50, line width=1pt, rounded corners=1.5pt]
    (0, -2.65) -- (\totalW - \tipBar, -2.65) -- (\totalW - \tipBar, -2.65 + \wingBar) -- (\totalW, -2.65 - \hBar/2) -- (\totalW - \tipBar, -2.65 - \hBar - \wingBar) -- (\totalW - \tipBar, -2.65 - \hBar) -- (0, -2.65 - \hBar) -- cycle;
\node[text=darkText, align=center] at (\totalW/2 - 0.2, -2.65 - \hBar/2) {
    \textbf{Duración acumulada:} 10 a 14 años de investigación y desarrollo
};

% --- Barra 2: Inversión requerida (> 1.000 millones de USD) ---
\filldraw[draw=grayText, fill=cardBorder!50, line width=1pt, rounded corners=1.5pt]
    (0, -3.75) -- (\totalW - \tipBar, -3.75) -- (\totalW - \tipBar, -3.75 + \wingBar) -- (\totalW, -3.75 - \hBar/2) -- (\totalW - \tipBar, -3.75 - \hBar - \wingBar) -- (\totalW - \tipBar, -3.75 - \hBar) -- (0, -3.75 - \hBar) -- cycle;
\node[text=darkText, align=center] at (\totalW/2 - 0.2, -3.75 - \hBar/2) {
    \textbf{Inversión acumulada:} Más de 1.000 millones de dólares (USD)
};

\end{tikzpicture}%
}
\caption[Pipeline tradicional de diseño y desarrollo de fármacos]{Pipeline tradicional de descubrimiento y desarrollo de fármacos, ilustrando la reducción progresiva de candidatos, tiempos e inversión requerida hasta la aprobación regulatoria. Adaptado de~\cite{chawla2024computational, hughes2011principles}.}
\label{fig:pipeline_tradicional}
\end{figure}
